\section{Introduction}
\label{sec:introduction}

Large language model deployment now depends on some form of adversarial
evaluation. Early red-teaming pipelines relied heavily on hand-authored probes,
but recent work has shown that language models themselves can generate diverse
test cases at scale \citep{perez2022redteaming,radharapu2023aart}. In parallel,
the jailbreak literature has made clear that alignment failures are not confined
to obscure edge cases: simple role-play prompts, adversarial suffixes, prompt
injection chains, and automated jailbreak generation can all elicit harmful or
policy-violating behavior from deployed systems
\citep{wei2023jailbroken,zou2023universal,yu2024dontlisten}. Evaluation work
has therefore shifted from merely finding attacks to measuring them with
standardized, reproducible metrics \citep{mazeika2024harmbench,souly2024strongreject}.

That shift creates a practical systems problem. In real evaluations, the hardest
part is often not writing another jailbreak prompt. It is building an execution
stack that can speak to multiple providers, preserve trajectories, keep track of
which failures are model refusals versus infrastructure errors, and emit outputs
that someone can actually act on. A paper-quality plot that ignores quota
exhaustion, invalid credentials, or stale registry entries may still look clean,
but it describes a world different from the one a red-team pipeline actually
inhabits.

This paper studies \system, a multi-provider red-teaming framework organized
around four ideas:
\begin{itemize}
\item a provider-agnostic execution layer with JSONL trajectory logging,
\item an ATT\&CK-inspired LLM attack taxonomy and campaign runner,
\item evaluator-integrated analysis and reporting, and
\item a new vulnerability engine centered on recon, focused attack, and
confirmation phases.
\end{itemize}

We do not present an idealized benchmark run. Instead, we write from the
archived evidence currently present in the repository. The main empirical basis
is a completed \texttt{compare --campaign full --evaluate strongreject} run
covering seven models with 20 probes each. A companion operational status table
records the remaining 13 configured aliases that were blocked by quota
exhaustion, invalid API credentials, or a stale Gemini preview identifier. This
produces a narrower but more honest result set than a typical leaderboard.

The paper makes four contributions. First, it documents the architecture of a
practical multi-provider red-teaming framework rather than a single benchmark.
Second, it describes a dedicated vulnerability engine that converts model traces
into confirmed findings and disclosure-ready outputs. Third, it reports a live
artifact-backed comparison over the reachable subset of the current registry,
including qualitative jailbreak traces on Mistral models. Fourth, it treats
operational reachability as a first-class evaluation result rather than a
background inconvenience.

The resulting claims are intentionally modest. The executed sample is small,
concentrated in the Claude and Mistral families, and judged by a single
\sr\ evaluator model. We therefore do not claim a broad cross-provider scaling
law. What we do claim is narrower and, we think, useful: a working framework
for live red teaming should report both model vulnerabilities and the execution
conditions that determine which models can be evaluated at all.
